\documentclass[../main.tex]{subfiles}
\begin{document}

% --------------------------------------------------
\section{layout.sty}
% --------------------------------------------------

\texttt{layout.sty}は, 紙面全体の見た目を調整するためのスタイルファイルである. 
背景色, 本文色, ヘッダー, フッター, 目次, タイトル, 見出し, 箇条書きの設定をまとめている. 

\subsection{背景色と本文色}

ページの背景色は\verb|\BaseColor|, 本文色は\verb|\TextColor|に設定している. 
これらの色名は\texttt{core.sty}で定義されているので, 文書全体の配色を変えたい場合は\texttt{core.sty}側の色コマンドを変更すればよい. 

\subsection{ヘッダーとフッター}

ページスタイル\texttt{note}を定義し, \verb|\pagestyle{note}|で文書全体に適用している. 
ヘッダーとページ番号の文字色には\verb|\AccentColor|を使う. 

\texttt{jlreq}クラスを\texttt{report}オプション付きで使う場合は, ヘッダーに章見出しと節見出しを表示する. 
それ以外の場合は, 節見出しを中心に表示する. 
どちらの場合もページ番号は小口側の下に置かれる. 

また, 章の開始ページなどで使われる\texttt{plain}ページスタイルについても, ページ番号の位置と色を揃えている. 
そのため, 通常ページと章開始ページでノンブルの位置がずれにくい. 

\subsection{目次}

目次の項目幅を調整し, 章や節の番号が本文と自然に揃うようにしている. 
また, \verb|\tableofcontents|を再定義し, 目次ページではヘッダーとフッターを出さない. 

本文中の見出しとPDFのリンク先がずれないように, \texttt{hyperref}が読み込まれている場合は目次登録の直前に\verb|\phantomsection|を入れる設定も行っている. 

\subsection{タイトルと概要}

\verb|\maketitle|を再定義している. 
\texttt{jlreq}クラスに\texttt{titlepage}オプションを付けた場合は, 独立したタイトルページを作る. 
付けない場合は, 本文の先頭にコンパクトなタイトルを出す. 

\texttt{abstract}環境も同じ方針で調整している. 
\texttt{titlepage}オプション付きの場合は概要を独立ページの中央に配置し, 通常の場合は本文冒頭の短いブロックとして表示する. 

\subsection{見出し}

章, 節, 小節の見出しは\texttt{jlreq}の\verb|\ModifyHeading|を使って調整している. 
章見出しは\texttt{report}オプション付きの場合だけ定義し, 中央揃えで大きく表示する. 

節見出しには左右に星印と罫線を入れて, ノートの区切りが視覚的に分かるようにしている. 
小節見出しは星印を添えた簡潔な形にしている. 

\subsection{箇条書き}

\texttt{enumitem}を読み込み, 番号付き箇条書きの表示を調整している. 
第1階層は\texttt{(1)}, 第2階層は\texttt{(1-1)}のように表示される. 
また, 本文との余白が広がりすぎないように, ラベルと本文の間隔や左マージンも設定している. 

\end{document}
